\clearpage
\VerdierPhysicalPageAnchor{0191}%
\begingroup\small\setlength{\parskip}{0pt}\setlength{\emergencystretch}{1.5em}
\setlength{\abovedisplayskip}{4pt plus 2pt minus 2pt}%
\setlength{\belowdisplayskip}{4pt plus 2pt minus 2pt}%

\VerdierSemanticAnchor{ch3d-2-2-2}
\noindent\textbf{Lemme 2.2.2.}\quad
\textit{Soit
$0\longrightarrow X\longrightarrow Y\longrightarrow Z\longrightarrow0$
une suite exacte de $\mathcal A$.}

\noindent\textit{a) Si $X$ et $Y$ sont des éléments de $M^{\wedge}$, il
existe un épimorphisme $u\colon L\longrightarrow Z$, $L\in M$.}

\noindent\textit{b) Si $X$ et $Z$ sont des éléments de $M^{\wedge}$, il
existe un épimorphisme $u\colon L\longrightarrow Y$, $L\in M$.}

\noindent\textit{c) Si $Y$ et $Z$ sont des éléments de $M^{\wedge}$, il
existe un épimorphisme $u\colon L\longrightarrow X$, $L\in M$.}

\VerdierSemanticAnchor{ch3d-2-2-3}
\medskip\noindent\textbf{Lemme 2.2.3.}\quad
\textit{Soit
$0\longrightarrow X\longrightarrow Y\longrightarrow Z\longrightarrow0$
une suite exacte de $\mathcal A$.}

\noindent\textit{a) Si $Y$ et $Z$ sont des éléments de $M^{\wedge}$, pour
tout épimorphisme $u\colon L\longrightarrow X$, $L\in M$, le noyau de
$u$ est un noyau d'un épimorphisme entre deux éléments de $M^{\wedge}$.}

\noindent\textit{b) Si $X$ et $Y$ sont des éléments de $M^{\wedge}$, pour
tout épimorphisme $u\colon L\longrightarrow Z$, $L\in M$, le noyau de
$u$ est un conoyau d'un monomorphisme entre deux éléments de
$M^{\wedge}$.}

\noindent\textit{c) Si $X$ et $Z$ sont des éléments de $M^{\wedge}$, pour
tout épimorphisme $u\colon L\longrightarrow Y$, $L\in M$, le noyau de
$u$ est un noyau d'un épimorphisme entre deux éléments de $M^{\wedge}$.}

\VerdierSemanticAnchor{ch3d-2-2-2-3-proof}
La démonstration des deux lemmes est laissée au lecteur. On remarquera que
les propriétés (a) et (c) des lemmes (2.2.2) et (2.2.3) ne font pas
intervenir la propriété $(E_2)$.

\VerdierSemanticAnchor{ch3d-2-2-4}
\medskip\noindent\textbf{Proposition 2.2.4.}\quad
\textit{a) L'ensemble $M^{\wedge}$ possède les propriétés $(E_6)$ et
$(E_2)$. L'ensemble $M$ est contenu dans l'ensemble $M^{\wedge}$. Tout
objet de $M^{\wedge}$ est quotient d'un objet de $M$.}

\noindent\textit{b) Tout ensemble d'objets $M'$ contenant $M$ tel que tout
objet de $M'$ soit un quotient d'un objet de $M$ et qui possède la
propriété $(E_3)$ est contenu dans $M^{\wedge}$.}

\noindent\textit{c) En particulier,
$(M^{\wedge})^{\wedge}=M^{\wedge}$.}

\VerdierSemanticAnchor{ch3d-2-2-4-proof}
La propriété a) résulte des deux lemmes précédents. Les propriétés b) et
c) sont immédiates.

\VerdierSemanticAnchor{ch3d-2-2-5}
\medskip\noindent\textbf{2.2.5.}\quad
Soient $M\subset M'$ deux ensembles d'objets de $\mathcal A$ tels que $M$
possède les propriétés $(E_4)$ et $(E_5)$, que $M'$ possède la propriété
$(E_4)$, et que tout objet de $M'$ soit quotient d'un objet de $M$. On dit
qu'un objet $X$ de $M'$ est de $M$-dimension inférieure ou égale à $n$ et
on note $M\text{-}\dim(X)\leq n$, s'il existe une suite exacte :
% Anomalie diplomatique de l'autorité : « inférieure où égale » ; correction : « inférieure ou égale ».

\VerdierSemanticAnchor{ch3d-2-2-5-exact-sequence}
\[
0\longrightarrow L_n\longrightarrow L_{n-1}\longrightarrow\cdots
\longrightarrow L_0\longrightarrow X\longrightarrow0,
\qquad L_i\in M,\quad0\leq i\leq n.
\]
\endgroup
